% chapters/42_safety.tex
% Part of: AI / LLM / Agents / Hardware Notes

\section{AI Safety Landscape \& Threat Models}

\subsection{Alignment problem overview}

% TODO

\subsection{Reward hacking \& Goodhart's law}

% TODO

\subsection{Deceptive alignment}

% TODO

\subsection{Containment \& control}

\textbf{The in-agent safety problem.} Current agent deployments rely heavily on
\emph{in-context} safety mechanisms: system prompts that say ``don't do X'',
tool descriptions with usage restrictions, and behavioural alignment instilled through
RLHF/DPO. These are soft controls --- they live inside the model's reasoning and can
fail under:
\begin{itemize}
  \item \textbf{Prompt injection:} malicious instructions embedded in observed content
    (web pages, emails, file contents) that hijack the agent's action stream.
  \item \textbf{Context overflow:} safety instructions scrolling out of the effective
    attention range in long sessions.
  \item \textbf{Reasoning drift:} the model convincing itself an exception is justified
    through multi-step chain-of-thought.
\end{itemize}

\textbf{Out-of-process policy enforcement} is the architectural response. Rather than
relying solely on the model's internal reasoning to stay within bounds, a runtime layer
\emph{outside} the agent intercepts actions before they execute and enforces hard
policy at the OS or network level. The model cannot override a kernel-level syscall
filter by reasoning; it simply receives a blocked-action error.

Real-world implementations:
\begin{itemize}
  \item \textbf{OpenShell (NVIDIA/NemoClaw):} policy engine + sandbox manager +
    privacy router. Uses Landlock, seccomp, and network namespaces. Routing decisions
    are based on operator policy files, not agent intent.
  \item \textbf{Cowork (Anthropic):} per-session Ubuntu VM with scoped workspace
    mount, product-level prohibited-action lists, and \texttt{permissionMode} tiers in
    the Claude Code agent runtime.
  \item \textbf{Browser security model analogy:} both approaches mirror how browsers
    sandbox tabs --- each session is isolated, cross-origin requests are blocked by
    default, and permissions (camera, clipboard, filesystem) require explicit grants.
\end{itemize}

\textbf{Principle of minimal footprint.} Anthropic's stated design principle for
agentic systems: prefer reversible over irreversible actions, request only necessary
permissions, avoid storing sensitive information beyond immediate need, and confirm
with users when scope is unclear. This is the soft-control complement to hard runtime
enforcement.

The emerging consensus is that robust agent containment requires \emph{both}: aligned
model behaviour (soft) to handle the common case, plus an out-of-process runtime
layer (hard) as the enforcement backstop when soft controls are insufficient.

\subsection{Training-time safety: data exclusion \& perplexity verification}

A complementary safety strategy operates at training time rather than at inference
time: deliberately \textbf{excluding dangerous data from the training corpus} so the
model never acquires the capability in the first place.

\textbf{Rationale.} Runtime filters and RLHF refusal training can suppress a
capability the model already has, but suppression can be bypassed via jailbreaks or
fine-tuning. If the dangerous capability was never learned, there is nothing to
unsuppress.

\textbf{Examples in practice.}
\begin{itemize}
  \item \textbf{EVO\,2 (biological foundation model):} all eukaryotic virus sequences
    were excluded from the 9-trillion-base-pair training dataset, including pathogens
    on the USDA Select Agents and Toxins list. The model was therefore never trained
    to model the statistical patterns of dangerous viruses. See Chapter~51.
  \item \textbf{CSAM prevention:} image generation models routinely filter known CSAM
    hashes from training data as a non-negotiable pre-processing step.
  \item \textbf{Weapons synthesis:} curated pretraining data pipelines for chemistry
    LLMs may exclude detailed synthesis routes for scheduled substances.
\end{itemize}

\textbf{Verifying exclusion works: perplexity probing.}
Removing data from a training corpus does not guarantee the model has no residual
representation of the excluded domain (e.g.\ from related sequences or indirect
references). \textbf{Perplexity probing} is the standard verification technique:

\begin{enumerate}
  \item Collect a held-out probe set of examples from the excluded domain (e.g.\ viral
    genome sequences).
  \item Feed these to the trained model and measure per-token perplexity.
  \item \textbf{High perplexity} $\Rightarrow$ the model has no internal model of these
    sequences; exclusion was effective.
  \item \textbf{Low perplexity} $\Rightarrow$ the model has acquired representations of
    the excluded domain via indirect exposure; exclusion was incomplete.
\end{enumerate}

For EVO\,2, viral sequence inputs produced high perplexity, and generation attempts
produced biological gibberish --- confirming successful exclusion.

\begin{notebox}
  Perplexity probing is the dual of perplexity-based contamination detection in eval
  methodology (Chapter~45): there, \emph{low} perplexity on benchmark examples reveals
  the model \emph{has} seen them. Here, \emph{high} perplexity on excluded data
  confirms the model \emph{has not}. Same technique, opposite direction.
\end{notebox}

\begin{warnbox}
  Data exclusion is necessary but not sufficient. Releasing open-source training code
  alongside a virus-excluded dataset still allows a malicious actor to obtain the
  excluded data independently and fine-tune their own version of the model. The EVO\,2
  authors acknowledge this dual-use risk explicitly.
\end{warnbox}

\textbf{Idempotency for agents:}
An action is \emph{idempotent} if executing it twice produces the same result as
executing it once. Critical when agents retry failed actions on real-world systems:
\begin{itemize}
  \item Non-idempotent retries produce duplicate side-effects: two emails sent, two
    orders placed, two support tickets created.
  \item Mitigation: pass an \emph{idempotency key} (unique ID per logical action) so
    the backend can detect and deduplicate retries. Standard practice in payment
    APIs (Stripe) and REST design (\texttt{PUT}/\texttt{DELETE} are idempotent;
    \texttt{POST} is not).
  \item Agent design rule: tools that mutate real-world state should be idempotent
    by design \emph{or} require explicit human confirmation and must never be retried
    automatically without an idempotency key.
\end{itemize}

